%--------------------
% Packages
% -------------------
\documentclass[11pt,a4paper]{article}
\usepackage[utf8x]{inputenc}
\usepackage[T1]{fontenc}
%\usepackage{gentium}
\usepackage{mathptmx} % Use Times Font
\usepackage{amsfonts}

\usepackage[pdftex]{graphicx} % Required for including pictures
\usepackage[swedish]{babel} % Swedish translations
\usepackage[pdftex,linkcolor=black,pdfborder={0 0 0}]{hyperref} % Format links for pdf
\usepackage{calc} % To reset the counter in the document after title page
\usepackage{enumitem} % Includes lists

\frenchspacing % No double spacing between sentences
\linespread{1.2} % Set linespace
\usepackage[a4paper, lmargin=0.1666\paperwidth, rmargin=0.1666\paperwidth, tmargin=0.1111\paperheight, bmargin=0.1111\paperheight]{geometry} %margins
%\usepackage{parskip}

\usepackage[all]{nowidow} % Tries to remove widows
\usepackage[protrusion=true,expansion=true]{microtype} % Improves typography, load after fontpackage is selected

\usepackage{lipsum} % Used for inserting dummy 'Lorem ipsum' text into the template


%-----------------------
% Set pdf information and add title, fill in the fields
%-----------------------
\hypersetup{
pdfsubject = {},
pdftitle = {},
pdfauthor = {}
}

\title{Direct neighbour basis for local dipole model}

%-----------------------
% Begin document
%-----------------------
\begin{document}
\maketitle
\section{Introduction}
To use a dipole model to perform MAD to bias IR spectra, we require the model to be differentiable w.r.t. atomic position. An existing model, derived from the tensor contraction of the virial tensor from rank-2 to rank-1, accurately predicts dipoles and spectra that agree with experiment using SOAP turbo descriptor inputs.

\begin{equation}
    \vec{\mu} = \sum_i \frac{\partial{U(q_i)}}{\partial{\vec r_{i}}}
    \label{eq:derivative-model}
\end{equation}

However, this model uses the first derivative of the SOAP turbo descriptors as its basis, leading to the requirement of second derivatives for use in MAD. This is undesirable, and an alternate model is pursued to eliminate the need for second derivatives of the SOAP turbo descriptors.

\section{Proposed Model}
We begin with the assumption that the total dipole is a sum of local dipole contributions as a function of local atomic environment. We choose to represent this local environment through SOAP turbo descriptors.

\begin{equation}
    \vec{\mu} = \sum_i \vec{\mu}_i
\end{equation}

From here, we propose that a sufficient ansatz is the sum of neighbour contributions of a scalar charge model output and a screened radial function, using the relative unit vector as a basis.

\begin{equation}
    \vec{\mu}_i = \sum_{j \in \mathcal{N}(i)} U(q_i)\, \hat{r}_{ij}\, f(r_{ij})
    \label{eq:local}
\end{equation}

Where $\vec{r}_{ij} = \vec{r}_i - \vec{r}_j$ is the relative position vector between centre atom $i$ and neighbour $j$, $r_{ij}$ its magnitude, $\hat{r}_{ij} = \vec{r}_{ij}/r_{ij}$ its unit vector, and $\mathcal{N}(i)$ the neighbours within $r_{\mathrm{cut}}$. $f(r_{ij})$ is the radial function containing the cutoff function $f_c$.

\begin{equation}
    f(r_{ij}) = r_{ij}^{\zeta}\, f_c(r_{ij}),
    \qquad
    f_c(r) = 1 - 3\left(\frac{r}{r_{\mathrm{cut}}}\right)^2 + 2\left(\frac{r}{r_{\mathrm{cut}}}\right)^3
\end{equation}

It is convenient to absorb the unit vector into the radial term. Writing
$\vec{r}_{ij} = r_{ij}\hat{r}_{ij}$ and
$R(r) = f(r)/r$, Eq.~\ref{eq:local} becomes

\begin{equation}
    \vec{\mu} = \sum_i \sum_{j \in \mathcal{N}(i)} U(q_i)\, R(r_{ij})\, \vec{r}_{ij},
    \qquad R(r) = r^{\zeta-1} f_c(r)
    \label{eq:disp}
\end{equation}

so that the screened functional is the single radial channel $r^{\zeta-1}f_c(r)$.
All vector character is carried by $\vec{r}_{ij}$; every other factor is an
invariant scalar.

\subsection{Initial Model Performance}

The model was implemented using a single hidden-layer neural network, optimised with the Separable Natural Evolution Strategy, in the same manner as the original Eq.~\ref{eq:derivative-model} model. The model underperforms on a liquid water dataset compared to the original, with the RMSE being in the neighbourhood of 6x the original derivative model. We speculate that this is not an issue with the model capacity, but with the ansatz of the radial term.

\section{Learned Radial Expansion}

We propose that we expand the radial term as a set of basis functions to allow the model to better resolve how each neighbour contribution relates to its position. The basis coefficients are included as parameters to be learned, and are extended to be indexed by (centre species, neighbour species) symmetric pairs. Doing so introduces a new discrete hyperparameter, $N_r$, that counts the number of radial basis functions.

A natural choice for the radial basis is to use Chebyshev polynomials of the
first kind:

\begin{equation}
    R_n(r_{ij}) = T_n(x_{ij})\, f_c(r_{ij}),
    \qquad
    x_{ij} = \frac{2 r_{ij}}{r_{\mathrm{cut}}} - 1,
    \qquad n = 0 \ldots N_r - 1
    \label{eq:cheb}
\end{equation}

where the polynomials are generated by the recurrence

\begin{equation}
    T_0(x) = 1,
    \qquad T_1(x) = x,
    \qquad T_{n+1}(x) = 2x\,T_n(x) - T_{n-1}(x),
    \label{eq:cheb-rec}
\end{equation}

evaluated once per pair and shared across all $N_r$ channels.

\subsection{Cutoff Function}

The choice of $f_c$ is important because the
smoothness of $\vec{\mu}$ at $r_{\mathrm{cut}}$ is a property of the cutoff. We
use the cosine cutoff which smoothly tends to zero at $r = r_{\mathrm{cut}}$.

\begin{equation}
    f_c(r) = \frac{1}{2}\left(1 + \cos \frac{\pi r}{r_{\mathrm{cut}}}\right),
    \qquad 0 \le r \le r_{\mathrm{cut}},
    \label{eq:fc}
\end{equation}

\subsection{Learned Coefficients and Species Pairs}

Each channel is given its own site scalar, so the network output is widened
from a scalar to a vector $U_i = (U_i^0 \ldots U_i^{N_r-1})$, with
$U_i^n = U^n(q_i)$ the $n$-th output for environment $q_i$. Each channel
carries a coefficient indexed by the unordered species pair $p(t_i, t_j)$ of
the bond:

\begin{equation}
    \vec{\mu} = \sum_i \sum_{j \in \mathcal{N}(i)} \sum_{n=0}^{N_r-1}
        c_n\!\left[p(t_i, t_j)\right] U_i^n\, R_n(r_{ij})\, \vec{r}_{ij}
    \label{eq:model}
\end{equation}

with $c \in \mathbb{R}^{N_r \times T(T+1)/2}$ for $T$ species, trained alongside the site parameters.

\section{Kernel Model}

The form of Eq.~\ref{eq:model} is linear in the site scalars $U_i^n$. If $U$ is
made linear in its parameters, the whole model becomes linear and the fit
reduces to a single ridge regression. This is the Gaussian approximation
potential (GAP) form of the model. The model, radial basis and cutoff are
unchanged; only the site function is replaced.

\subsection{Kernel and Sparse Points}

Similarity between two environments is measured by the polynomial
dot-product kernel on the unit-normalised descriptors,

\begin{equation}
    k(q_i, q_m) = (q_i \cdot q_m)^{\xi}, \qquad \xi \in \mathbb{N},
    \label{eq:kernel}
\end{equation}

where $\xi$ is a hyperparameter that sharpens the kernel. For each species
$t$ a set $\mathcal{M}_t$ of $M$ representative environments is drawn from
the training set by farthest-point sampling. An environment of species $t$ is compared only with the sparse
points of species $t$.

\subsection{Site Scalars and Species Pairs}

Each channel of the site scalar is a kernel expansion over the sparse set of
the centre species,

\begin{equation}
    U_i^n = \sum_{m \in \mathcal{M}_{t_i}} \alpha_m^{n}\, k(q_i, q_m).
    \label{eq:gap-site}
\end{equation}

In Eq.~\ref{eq:model} the pair coefficient $c_n[p(t_i,t_j)]$ multiplies
$U_i^n$, so the product $c_n \alpha_m^n$ is one free weight per (sparse
point, channel, neighbour species). It is therefore absorbed into $\alpha$,
which gains a neighbour-species index $t'$. Collecting the geometry of each
centre by neighbour species,

\begin{equation}
    \vec{v}_i^{\,n,t'} = \sum_{j \in \mathcal{N}(i),\, t_j = t'}
        R_n(r_{ij})\, \vec{r}_{ij},
    \label{eq:geom}
\end{equation}

the model becomes

\begin{equation}
    \vec{\mu} = \sum_i \sum_{m \in \mathcal{M}_{t_i}} \sum_{n=0}^{N_r-1}
        \sum_{t'} \alpha_m^{n,t'}\, k(q_i, q_m)\, \vec{v}_i^{\,n,t'}.
    \label{eq:gap}
\end{equation}

\subsection{Linear Fit}

Eq.~\ref{eq:gap} is linear in $\alpha$. For a training frame $f$ the design
block has one three-component column per weight,

\begin{equation}
    \Phi_f\big[:, (m,n,t')\big] = \sum_{i \in f,\; t_i = t(m)}
        k(q_i, q_m)\, \vec{v}_i^{\,n,t'} \;\in\; \mathbb{R}^3,
\end{equation}

and stacking $F$ frames gives $\Phi \in \mathbb{R}^{3F \times N_{\mathrm{col}}}$ with
$N_{\mathrm{col}} = N_r\, T \sum_t M_t$. The weights are the ridge solution

\begin{equation}
    \alpha = \arg\min_\alpha \|\Phi\alpha - \vec{\mu}^{*}\|^2
        + \lambda_{\mathrm{eff}} \|\alpha\|^2,
    \qquad
    \lambda_{\mathrm{eff}} = \lambda\, \overline{\|\Phi_{:,c}\|^2},
    \label{eq:ridge}
\end{equation}

\begin{equation}
    \alpha = \Phi^{\mathsf{T}} \left(\Phi\Phi^{\mathsf{T}}
        + \lambda_{\mathrm{eff}} I\right)^{-1} \vec{\mu}^{*},
    \label{eq:dual}
\end{equation}

\subsection{Model Performance}

Test RMSE on 500 held-out frames of the liquid water dataset, in
e$\cdot$bohr.

\begin{center}
\begin{tabular}{lr}
\hline
Model & RMSE \\
\hline
ANN, best original Eq.~\ref{eq:derivative-model} model & 0.0847 \\
ANN, initial Eq.~\ref{eq:local} screened, $\zeta = -2$ & 0.4949 \\
ANN, Chebyshev species-learned, Adam optimiser, $N_r = 10$ & 0.1373 \\
ANN, Chebyshev species-learned, Adam optimiser, $N_r = 20$ & 0.1137 \\
GAP, $M = 2000$, $\zeta = 3$, $N_r = 22$ & 0.1485 \\
\hline
\end{tabular}
\end{center}

\section{Towards an Improved Model}\label{sec:scope}

From here, accuracy will be quoted as relative RMSE (RRMSE)
\begin{equation}
\mathrm{RRMSE} = \frac{\mathrm{RMSE}}{\mathrm{std}(\mathrm{reference})},\qquad R^2 = 1-\mathrm{RRMSE}^2 ,
\label{eq:rrmse}
\end{equation}
computed over components (the three Cartesian dipole components; the three diagonal and the three
off-diagonal components of the polarisability tensor separately). In the GAP fits below the
regularisation is expressed through the observation-noise parameter $\sigma$ of the GAP
estimator, which plays the role of $\lambda_{\mathrm{eff}}$ in Eq.~\ref{eq:ridge}: the squared
residual of each frame is weighted by $1/(\sigma^2 N)$ with $N$ its number of atoms, so a larger
$\sigma$ regularises more strongly.

Published reference values from the TNEP paper (Xu et al. 2024): dipole 0.85\,\%; polarisability 16.3\,\% (diagonal) / 20.4\,\%
(off-diagonal), on the susceptibility $\chi/\rho$.

\section{The Expanded Model Family}\label{sec:readout}

We have good evidence from the previously tested dipole models that our species-learned basis expansion performs well, so we generalise the formulation to include polarisability as a learned tensor. To keep $\alpha$ for the kernel weights, the polarisability tensor is written $\mathbf{A}$. Both properties are sums over centres of site scalars multiplying fixed geometric factors, a vector $\vec v_i^{\,k}$ for the dipole and a symmetric $3\times3$ tensor $\mathbf{V}_i^{\,k}$ for the polarisability, where the index $k$ runs over radial channels, neighbour species and, for the polarisability, kernel groups:

\begin{equation}
\vec{\mu} = \sum_i \sum_k U^{k}_i\, \vec v_i^{\,k}, \qquad
\mathbf{A} = \sum_i \sum_k U^{k}_i\, \mathbf{V}_i^{\,k} .
\label{eq:readout}
\end{equation}

Each site scalar $U_i^k$ is a kernel expansion of the form of Eq.~\ref{eq:gap-site}, so every model in this family is linear in its weights and is fitted as in Eq.~\ref{eq:ridge}. The radial channels are the warped Chebyshev family of the production readout,

\begin{equation}
R_n(r) = T_n(x)\,f_c(r), \qquad x = 2\left(\frac{r}{r_{\mathrm{cut}}}-1\right)^2-1,
\label{eq:cheb-warped}
\end{equation}

which replaces the linear map of Eq.~\ref{eq:cheb}, with a cosine (Eq.~\ref{eq:fc}) or squared-cosine ($f_c^2$) cutoff, $N_r = 10$ channels, and independent coefficients per (channel, neighbour species) pair as in Eq.~\ref{eq:gap}.

\section{Dipole models}\label{sec:dipole}

\subsection{The Established Model}\label{sec:dip-original}

Eq.~\ref{eq:gap} with $N_r=10$ warped Chebyshev channels per neighbour species on $r_{\mathrm{cut}}=3.0$~\AA{} with a
squared-cosine cutoff, each channel with its own site function: a sparse kernel on the compressed
$l_{\max}=n_{\max}=4$ SOAP (hard/soft cutoff 3.5/3.0~\AA{}), $M=250$ farthest-point
environments per species, $\xi=2$, $\sigma=10^{-3}$. $N_{\mathrm{col}} = 10\,000$. \textbf{Test RRMSE 1.76\,\%}.

\subsection{Model Improvements}\label{sec:dip-local}

Two improvements can be made to the already-established dipole:

\begin{enumerate}[leftmargin=*,nosep]
\item \textbf{A longer, uncompressed descriptor.} $l_{\max}=n_{\max}=7$ with no compression (the full
  840-dimensional power spectrum). 1.76 $\rightarrow$
  \textbf{1.56\,\%}.
\item \textbf{An additional two-body kernel.} A second GAP which uses a kernel of the radial distribution of the
  environment on the same 3~\AA{}, its own FPS set ($M=250$) and its own prior,
  $\delta_{2b}=0.1$. 1.56 $\rightarrow$ \textbf{1.49\,\%} at
  $\sigma=3\times10^{-3}$. $N_{\mathrm{col}}$ doubles to 20\,000.
\end{enumerate}

\subsection{The Additional Two-body GAP}\label{sec:2b}

The established model uses a single sparse kernel on the SOAP vector $q_i$. The
NEP decomposition of the site energy by body order suggests adding a second kernel on the
radial distribution of the environment. Collecting the radial channels of Eq.~\ref{eq:geom}
without their bond vectors gives, per centre, the invariant vector

\begin{equation}
q^{(2)}_{i,\,n t'} = \sum_{j \in \mathcal{N}(i),\, t_j = t'} R_n(r_{ij}),
\qquad n = 0 \ldots N_r-1,
\label{eq:2b-desc}
\end{equation}

($N_r T$ numbers), which is unit-normalised and compared with the same kernel as
Eq.~\ref{eq:kernel}, $k^{(2)}(q^{(2)}_i, q^{(2)}_{m}) = (q^{(2)}_i \cdot q^{(2)}_{m})^{\xi}$, over its
own farthest-point set $\mathcal{M}^{(2)}_t$ ($M=250$ per species). Each channel's site scalar becomes the sum of the two expansions,

\begin{equation}
U_i^{n,t'} = \sum_{m \in \mathcal{M}_{t_i}} \alpha_m^{n,t'}\, k(q_i, q_m)
\;+\;\delta_{2b} \sum_{m' \in \mathcal{M}^{(2)}_{t_i}} \beta_{m'}^{n,t'}\, k^{(2)}(q^{(2)}_i, q^{(2)}_{m'}),
\label{eq:2b-site}
\end{equation}

with its own weights $\beta$ and its own prior; $\delta_{2b}$ is the relative prior amplitude of the
two-body kernel. The geometry $\vec v_i^{\,n,t'}$ of Eq.~\ref{eq:geom} is unchanged.

\subsection{Improved Dipole Results}\label{sec:dip-results}

Test set of 500 frames; RMSE per Cartesian component in e$\cdot$bohr per cell.

\begin{center}
\begin{tabular}{p{7.6cm}rr}
\hline
Model & RRMSE & RMSE \\
\hline
ANN, best original Eq.~\ref{eq:derivative-model} model & 1.36\,\% & 0.0847 \\ % RRMSE derived as 0.0847 / 6.24 (test spread per component); verify same test set and definition
Original GAP, $l_{\max}=n_{\max}=4$ compressed SOAP & 1.76\,\% & 0.110 \\
GAP, $l_{\max}=n_{\max}=7$ uncompressed SOAP & 1.56\,\% & 0.097 \\
GAP, $l_{\max}=n_{\max}=7$ uncompressed SOAP + two-body GAP, $\delta_{2b}=0.1$ & 1.49\,\% & 0.093 \\
Published TNEP (NEP descriptors, 200k SNES generations) & 0.85\,\% & 0.053 \\
\hline
\end{tabular}
\end{center}

\section{Polarisability Models}\label{sec:pol}

\subsection{Tensor Readout}\label{sec:pol-readout}

The polarisability model is Eq.~\ref{eq:readout} with symmetric $3\times3$ tensors
$\mathbf{V}_i^{\,k}$ in place of the vectors $\vec v_i^{\,k}$. Each tensor is built from the
bond vectors $\vec r_{ij}$ of the neighbours of $i$ with the radial channels of
Eq.~\ref{eq:cheb-warped}, so that, as for the dipole, all tensor character is carried by the
geometry and every other factor is an invariant scalar. The site scalars are grouped by the
kind of tensor they multiply; each group has its own kernel expansion of the form of
Eq.~\ref{eq:gap-site}. Three groups are used. Throughout, $N_r = 10$ channels on
$r_{\mathrm{cut}} = 4$~\AA{} with the cosine cutoff, and $\mathbf{I}$ is the $3\times3$ identity.

\subsection{Isotropic Bond Group}\label{sec:pol-iso}

The first group carries the trace of the tensor. For each channel and neighbour species,

\begin{equation}
\mathbf{V}_i^{\,\mathrm{iso},\,n,t'} = \sum_{j \in \mathcal{N}(i),\, t_j = t'} R_n(r_{ij})\, r_{ij}^2\, \mathbf{I},
\label{eq:pol-iso}
\end{equation}

a multiple of the identity, which can only move the three diagonal components, and by the same
amount. Its site scalars are written $U_i^{\,\mathrm{iso},\,n,t'}$.

\subsection{Traceless Bond Group}\label{sec:pol-tr}

The second group carries the anisotropy of each bond,

\begin{equation}
\mathbf{V}_i^{\,\mathrm{tr},\,n,t'} = \sum_{j \in \mathcal{N}(i),\, t_j = t'} R_n(r_{ij})
\left(\vec r_{ij}\,\vec r_{ij}^{\,\mathsf{T}} - \frac{r_{ij}^2}{3}\,\mathbf{I}\right),
\label{eq:pol-tr}
\end{equation}

which has zero trace by construction, so it can only move the shape of the tensor. Its site
scalars are $U_i^{\,\mathrm{tr},\,n,t'}$, independent of those of the isotropic group. The
isotropic and traceless groups together are the \textbf{original model}:
$l_{\max}=n_{\max}=7$ compressed SOAP, $M=250$, $\sigma=0.1$, $N_{\mathrm{col}} = 20\,000$;
\textbf{16.6 / 19.9\,\%} diagonal / off-diagonal on 500 training frames (16.9 / 20.2 on 400),
0.63~ms/atom. Every later number is measured against it.

\subsection{Cross Group}\label{sec:pol-cross}

The third group is built from the dipole geometry vectors of Eq.~\ref{eq:geom}. Taking the first
four radial channels, the vectors $\vec v_i^{\,a}$ with $a$ running over these (channel,
neighbour species) combinations are combined pairwise into symmetrised, traceless products,

\begin{equation}
\mathbf{V}_i^{\,\times,\,ab} = \frac{1}{2}\left(\vec v_i^{\,a}\,\vec v_i^{\,b\,\mathsf{T}} + \vec v_i^{\,b}\,\vec v_i^{\,a\,\mathsf{T}}\right)
- \frac{1}{3}\left(\vec v_i^{\,a}\cdot\vec v_i^{\,b}\right)\mathbf{I},
\qquad a \le b,
\label{eq:pol-cross}
\end{equation}

with site scalars $U_i^{\,\times,\,ab}$. Because each $\vec v_i^{\,a}$ is a sum over neighbours,
the product contains terms $\vec r_{ij}\,\vec r_{ik}^{\,\mathsf{T}}$ with two different bonds, which
the bond groups of Eqs.~\ref{eq:pol-iso} and \ref{eq:pol-tr} cannot represent.

\subsection{The Full Model and its Fit}\label{sec:pol-form}

Assembling the three groups,

\begin{equation}
\begin{array}{rl}
\mathbf{A} = & \displaystyle \sum_i \sum_{n,t'} \Big[\, U_i^{\,\mathrm{iso},\,n,t'}\, \mathbf{V}_i^{\,\mathrm{iso},\,n,t'}
\;+\; U_i^{\,\mathrm{tr},\,n,t'}\, \mathbf{V}_i^{\,\mathrm{tr},\,n,t'} \Big] \\[2ex]
 & \displaystyle +\; \sum_i \sum_{a \le b} U_i^{\,\times,\,ab}\, \mathbf{V}_i^{\,\times,\,ab} ,
\end{array}
\label{eq:pol-form}
\end{equation}

with every site scalar a kernel expansion as in Eq.~\ref{eq:gap-site}. The model is linear in
the weights, and the fit is that of Eq.~\ref{eq:ridge} with six rows per frame (the three diagonal
and three off-diagonal components of $\mathbf{A}$) in place of three, so that
$\Phi \in \mathbb{R}^{6F \times N_{\mathrm{col}}}$.

\subsection{Polarisability Results}\label{sec:pol-results}

Test set of 500 frames, fits on 500 training frames; RRMSE and RMSE given as diagonal /
off-diagonal, RMSE per component in \AA$^3$ per cell.

\begin{center}
\begin{tabular}{p{6.0cm}rrr}
\hline
Model & RRMSE & RMSE & ms/atom \\
\hline
Two bond groups, SNES-trained ANN (15 neurons, 25\,k parameters, $\sigma_0=0.02$, 20\,000 generations) & 35.7 / 27.5\,\% & 0.317 / 0.129 & \\
Two bond groups: the original model, Sec.~\ref{sec:pol-tr} & 16.6 / 19.9\,\% & 0.147 / 0.093 & 0.63 \\
+ cross group, four channels (Sec.~\ref{sec:pol-cross}) & 14.9 / 16.0\,\% & 0.132 / 0.075 & 0.86 \\
+ second SOAP scale at 6.5~\AA{}, $\sigma=0.3$ & 13.9 / 16.2\,\% & 0.123 / 0.076 & 2.5 \\
+ induced coupling + quadrupole + 6.5~\AA{} scale, $\sigma=0.3$ & 13.0 / 15.2\,\% & 0.115 / 0.071 & 2.6 \\
\quad same, $\sigma=0.1$ & 13.9 / 14.6\,\% & 0.123 / 0.068 & 2.6 \\
Published TNEP (same test frames; paper RMSE per molecule converted per component per cell) & 16.3 / 20.4\,\% & 0.148 / 0.101 & \\
\hline
\end{tabular}
\end{center}

\end{document}
